% Taller 2 · Filosofía 10.° · Periodo IV · «Diseña tu isla de Utopía» (semana 8, clase 1)
% Imagen: img/utopia.jpg (Holbein, 1518; Wikimedia Commons, dominio público).
% Desempeño: D5 (FIL-10-P4-D5). Una hora (≈55 min), en grupos de 3 o 4. Materiales: esta hoja,
% una hoja en blanco por grupo para el mapa y marcadores.
% Compilar: bash curriculo/plantilla/compilar.sh curriculo/filosofia/10/taller-2-P4.tex
\documentclass[taller]{planeacion}
\materia{filosofia}
\grado{10.°} \periodo{IV} \anio{2027}
\begin{document}
\fichataller{2}{Semana 8, clase 1}{55 min}{Grupos de 3–4}

\titulotaller{Diseña tu isla de Utopía}{En grupos de 3 o 4. Lo escrito va \textbf{en esta hoja}; el mapa, \textbf{en la hoja
  en blanco} que entrega el docente.}

\reto{0}{Lo que hizo Moro}{5 min}{0,5}
En el libro I, Moro critica a la Inglaterra de su tiempo, donde los campos de cultivo se volvían
pastos para ovejas y los campesinos quedaban en la miseria:
\fuente{Vuestras ovejas, que solían ser tan mansas y comer tan poco, ahora, según dicen, se han
  vuelto tan voraces y salvajes que se comen a los mismos hombres.}
  {Tomás Moro, \emph{Utopía} (1516), libro I (traducción adaptada)}
\imagenlado[0.3]{img/utopia.jpg}{En el libro II describe la isla de Utopía: allí no hay propiedad privada,
  todos trabajan solo seis horas al día y las familias cambian de casa cada diez años por sorteo.
  Moro critica un problema y luego imagina una solución radical: eso harán ustedes.\par\vspace{3pt}
  {\footnotesize \emph{Mapa de la isla de Utopía}, grabado de Ambrosius Holbein para la edición de
  \emph{Utopía} de 1518. Wikimedia Commons, dominio público.}}
\begin{tareas}
  \tarea{\envozalta respondan en el grupo: ¿qué critica Moro con la imagen de las ovejas?}
  \tarea{Observen el mapa: ¿cómo es la isla y qué la separa del resto del mundo? ¿Por qué creen
    que Moro la imaginó como una isla?}
\end{tareas}

\reto{0}{El problema}{10 min}{1,0}
\begin{tareas}
  \tarea{Escojan \textbf{un problema real} del municipio o del colegio (basuras, movilidad, uso del
    celular, convivencia, agua…).}
  \tarea{\enlahoja escríbanlo en una frase y pónganle nombre a su isla.}
\end{tareas}
\textbf{Nombre de la isla:}\lineas{1}
\textbf{El problema que queremos resolver (como las ovejas de Moro):}\lineas{2}

\reto{0}{La isla y sus leyes}{30 min}{2,5}
\begin{tareas}
  \tarea{\textbf{En la hoja en blanco,} dibujen el \textbf{mapa de su isla}.}
  \tarea{\enlahoja escriban \textbf{tres leyes} que resuelvan el problema, cada una con su razón.}
  \tarea{Para cada ley, escriban \textbf{un riesgo}: ¿qué podría salir mal o a quién podría perjudicar?}
\end{tareas}
\textbf{Ley 1, su razón y su riesgo:}\lineas{3}
\textbf{Ley 2, su razón y su riesgo:}\lineas{3}
\textbf{Ley 3, su razón y su riesgo:}\lineas{3}

\reto{0}{Exposición y cierre}{10 min}{1,0}
\begin{tareas}
  \tarea{\envozalta en un minuto por grupo: el problema, la mejor ley y su riesgo.}
  \tarea{\enlahoja respondan en dos líneas: ¿es posible una sociedad perfecta?}
\end{tareas}
\lineas{2}

\apartado{Cómo se valora}
\begin{rubrica}
  \criterio{Problema}{Real, concreto y bien explicado}{Real pero general}{Vago}{Ausente}
  \criterio{Leyes y razones}{Tres leyes que resuelven el problema, con razones}{Tres leyes, alguna sin razón}{Leyes sin relación clara con el problema}{Menos de dos leyes}
  \criterio{Riesgos}{Un riesgo serio y honesto por ley}{Riesgos para casi todas}{Riesgos superficiales}{Sin riesgos}
  \criterio{Exposición}{Clara, en un minuto, responde preguntas}{Clara}{Confusa}{No expone}
\end{rubrica}
\cierretaller{5,0}
\end{document}
